% !PRINTED 590
\weberpage{590}
%<<CONT
der $\varkappa$, von Einheitsfactoren abgesehen, mit der Gesammtheit
der $\pi$ übereinstimmen muss.

Es handelt sich nun noch darum, die Primzahlen des
Körpers $R(i)$ wirklich zu ermitteln.

Jede Primzahl ist Theiler von unendlich vielen rationalen
ganzen Zahlen, z. B. von der Norm und allen ihren Vielfachen.
Unter den rationalen positiven ganzen Zahlen, die durch eine Primzahl
 $\pi$ theilbar sind, ist eine, die wir mit $n$ bezeichnen wollen,
die kleinste, und diese muss eine Primzahl im Körper $R$ sein.
Denn wäre sie in $R$ zerlegbar, so müsste nach 5. einer ihrer Factoren,
also eine noch kleinere Zahl, durch $\pi$ theilbar sein. Ebenso
ist auch umgekehrt jede reelle Primzahl $n$ wenigstens durch eine
Primzahl $\pi$ theilbar. Es sei also $n = \pi\alpha$, woraus durch Bildung
der Norm $n^2 = N(\pi)N(\alpha)$ folgt. Da $n$ eine Primzahl ist und
$N(\pi)N(\alpha)$ ganze Zahlen, so sind zwei Fälle möglich:
\[
\begin{aligned}
1.\quad & N(\pi) = n, \quad N(\alpha) = n,\\
2.\quad & N(\pi) = n^2, \quad N(\alpha) = 1.
\end{aligned}
\]

Im ersten Falle ist, wenn $\pi'$ die zu $\pi$ conjugirte Zahl bedeutet
 und $\pi = a+bi$ gesetzt wird, $n = \pi\pi' = a^2+b^2$,
$\alpha = \pi'$, und man sieht, dass dieser Fall nur dann eintritt, wenn
$n$ zu den Primzahlen $p$ gehört, die in die Summe von zwei
Quadraten zerlegbar sind. Umgekehrt kann jede solche Zahl $p$
in zwei conjugirte Factoren $\pi\pi'$ zerlegt werden, deren keine
eine Einheit ist. \gesperrt{Die Primzahlen} $p$ \gesperrt{sind also im Körper}
$R(i)$ \gesperrt{nicht Primzahlen.}

Im zweiten Falle, der hiernach bei den Primzahlen $q$ eintritt,
ist $\alpha$ eine Einheit, also $n$ mit $\pi$ associirt, d. h. die \gesperrt{reellen
Primzahlen} $q$ \gesperrt{sind auch im Körper} $R(i)$ \gesperrt{Primzahlen.}

Die Primzahl 2 gehört, wie schon bemerkt, zu der ersten
Art, und es ist $2 = (1+i)(1-i)$. Aber sie nimmt eine besondere
 Stellung ein, weil $1-i = -i(1+i)$, also $2 = -i(1+i)^2$
ist. Die reelle Primzahl 2 ist also (von dem Factor $-i$ abgesehen)
im Körper $R(i)$ \gesperrt{das Quadrat einer Primzahl.} Dieser
Fall tritt, wie man leicht sieht, bei keiner der übrigen Primzahlen
 $p$ ein.

% [sic] German authority prints "des Körper R(i)" for grammatical "des Körpers R(i)"; diplomatic DE preserved.
Die Gesammtheit der Primzahlen des Körper $R(i)$ besteht
also aus den reellen Primzahlen $q$ und aus den Factoren der Primzahlen
 $p$. Ist $p$ in die Summe zweier Quadrate, $a^2+b^2$, zerlegt, so
kennt man auch die conjugirten Factoren $a\pm bi$ von $p$. Dazu
%CONT>>
